\documentclass[11pt]{article}

\usepackage[french]{babel}

\usepackage[T1]{fontenc}

\usepackage{amsfonts}
\usepackage{amsmath}
\usepackage{amssymb}

\usepackage{graphicx}
\usepackage{multirow}
\usepackage{multicol}
\usepackage{xcolor}
\usepackage[shortlabels]{enumitem}

\usepackage[margin=2cm]{geometry}

\usepackage{mathtools}
\usepackage{siunitx}
\usepackage{physics}

\usepackage{tikz}
\usepackage{hyperref}

\usepackage{algorithm}
\usepackage{algpseudocode}
\usepackage{listings}
\usetikzlibrary{quantikz2,babel}

\tikzset{qcircuit/.style={
    column sep=.9cm,
    row sep=0.7cm,
    operator/.style={draw=black!80, line width=0.5pt, fill=gray!5, minimum height=0.7cm, minimum width=0.7cm},
    phase/.style={draw=black!80, fill=black!20, circle, minimum size=4pt},
    slice/.style={font=\footnotesize\itshape, color=gray!50, dashed}
}}


\DeclarePairedDelimiter\ceil{\lceil}{\rceil}
\DeclarePairedDelimiter\floor{\lfloor}{\rfloor}


\AtBeginDocument{\RenewCommandCopy\qty\SI}

\makeatletter
   \def\vhrulefill#1{\leavevmode\leaders\hrule\@height#1\hfill \kern\z@}
\makeatother

\newcommand{\moy}[1]{\langle #1 \rangle}
\newcommand{\pDer}[2]{\frac{\partial #1}{\partial #2}}
\newcommand{\pDerd}[2]{\frac{\partial^2 #1}{\partial #2^2}}

\newcounter{questioncounter}
\newenvironment{question}[1]{
\refstepcounter{questioncounter}
\noindent\rule{0.05\textwidth}{1pt}\,\,\raisebox{-.3\ht\strutbox}{\textbf{Question \thequestioncounter \hspace{5mm} #1}}\,\,\vhrulefill{1pt}\par\vspace{0.4cm}
}{\par}

%======================================================
%en-tete
%======================================================
\begin{document}
\selectlanguage{french}
\noindent{\Large Anthony Dubeau: Duba0207 } \hfill {\Large \today}\\
{\Large Anthony Wroblewski: Wroa8915}


\par\vspace{5mm}
\begin{center}
   {\Large IFT436 - Algorithmes et structures de données} \\ \vspace{5mm}
   {\LARGE\sffamily Devoir \#2} \\ \vspace{5mm}
\end{center}

%======================================================
%question 1
%======================================================
\begin{question}{Une question d’échauffement}
\begin{enumerate}[label=(\alph*)]

\item On démontre les deux implications.

\medskip
\noindent
\(\Rightarrow\) Supposons que \(G\) est fortement connexe.
Pour tout sommet \(u\in V\) et tout sommet \(v\in V\), il existe
un chemin orienté de \(u\) à \(v\) dans \(G\). Ainsi,
\[
r(G,u)=V.
\]
Il existe également un chemin orienté de \(v\) à \(u\) dans \(G\).
En renversant chacune de ses arêtes, on obtient un chemin orienté
de \(u\) à \(v\) dans \(G^T\). Par conséquent,
\[
r(G^T,u)=V.
\]
Les deux égalités sont donc vérifiées pour tout \(u\in V\).

\medskip
\noindent
\(\Leftarrow\) Supposons maintenant que, pour tout \(u\in V\),
\[
r(G,u)=V
\qquad\text{et}\qquad
r(G^T,u)=V.
\]
Soient \(u,v\in V\) deux sommets quelconques.
Puisque \(v\in r(G,u)\), il existe un chemin orienté de \(u\)
à \(v\) dans \(G\).
De plus, puisque \(v\in r(G^T,u)\), il existe un chemin orienté
de \(u\) à \(v\) dans \(G^T\). En renversant ses arêtes,
on obtient un chemin orienté de \(v\) à \(u\) dans \(G\).

Il existe donc un chemin orienté dans les deux sens entre
toute paire de sommets. Ainsi, \(G\) est fortement connexe,
ce qui complète la démonstration.

\item 
Soit \(u\in V\) un sommet tel que
\[
r(G,u)=r(G^T,u)=V.
\]
Montrons d'abord que \(G\) est fortement connexe.
Soient \(v,w\in V\) deux sommets quelconques.

Puisque \(v\in r(G^T,u)\), il existe un chemin orienté
de \(u\) à \(v\) dans \(G^T\), donc un chemin orienté
de \(v\) à \(u\) dans \(G\).
De plus, puisque \(w\in r(G,u)\), il existe un chemin
orienté de \(u\) à \(w\) dans \(G\).

En passant par \(u\), on peut donc aller de \(v\) à \(w\).
Comme \(v\) et \(w\) sont quelconques, \(G\) est fortement connexe.
D'après le résultat de (a), on obtient alors
\[
\boxed{r(G,v)=r(G^T,v)=V
\qquad\text{pour tout }v\in V.}
\]
\item
On suppose que le graphe est représenté par des listes
d'adjacence. L'algorithme est le suivant :
\begin{enumerate}[label=\roman*]
    \item Choisir un sommet arbitraire \(u\in V\).
    \item Effectuer un parcours en largeur (BFS) dans \(G\)
    à partir de \(u\). Si tous les sommets ne sont pas visités,
    retourner \texttt{faux}.
    \item Construire le graphe transposé \(G^T\) en renversant
    chaque arête de \(G\).
    \item Effectuer un BFS dans \(G^T\) à partir du même sommet
    \(u\). Si tous les sommets sont visités, retourner
    \texttt{vrai}; sinon, retourner \texttt{faux}.
\end{enumerate}


Les deux parcours permettent de déterminer respectivement
\(r(G,u)\) et \(r(G^T,u)\). D'après (b), si ces deux ensembles
sont égaux à \(V\), alors \(G\) est fortement connexe.
Réciproquement, d'après (a), si \(G\) est fortement connexe,
les deux parcours visitent tous les sommets.
L'algorithme retourne donc \texttt{vrai} si et seulement si
\(G\) est fortement connexe.

\medskip
\textbf{Complexité.}
Comme vue en classe, Chaque BFS visite chaque sommet au plus une fois et examine
chaque arête au plus une fois. Il coûte donc
\(O(|V|+|E|)\).
La construction de \(G^T\) consiste à créer une liste
d'adjacence pour chaque sommet et à parcourir chaque arête
\((v,w)\) pour ajouter l'arête \((w,v)\). Elle coûte également
\(O(|V|+|E|)\).

La complexité totale est donc
\[
\boxed{O(|V|+|E|)}.
\]
Je pense que cette complexité est optimale dans le pire cas pour une représentation par listes d'adjacence car il peut être nécessaire
d'examiner tout les sommets et d'arêtes pour
décider de la forte connexité.
\item
L'idée est de trouver la composante fortement connexe contenant
un sommet \(u\) en effectuant deux parcours en largeur (BFS).

Un BFS dans \(G\) à partir de \(u\) trouve les sommets accessibles
depuis \(u\), soit \(r(G,u)\). Un BFS dans \(G^T\) à partir du même
sommet trouve les sommets qui peuvent atteindre \(u\) dans \(G\),
soit \(r(G^T,u)\). Leur intersection
\[
X=r(G,u)\cap r(G^T,u)
\]
contient donc exactement les sommets qui sont mutuellement
accessibles avec \(u\).

On enregistre cette composante, puis on recommence avec un sommet
qui n'appartient à aucune composante déjà trouvée. Le graphe
transposé est construit une seule fois au début.

\medskip
\textbf{Pseudo-code.}\\
On suppose que \(\operatorname{BFS}(G,u)\) retourne l'ensemble
des sommets visités.
% \begin{lstlisting}
% ComposantesFortementConnexes(G):

%     construire G_transpose en renversant chaque arete de G
%     composantes <- liste vide
%     non_classes <- V(G)

%     tant que non_classes est non vide:

%         choisir u dans non_classes

%         accessibles <- BFS(G, u)
%         vers_u <- BFS(G_transpose, u)
%         X <- intersection(accessibles, vers_u)

%         ajouter X a composantes
%         non_classes <- non_classes \ X

%     retourner composantes
% \end{lstlisting}

%pas content de comment il sortait sur le pdf fac je l'ai refait avec tabbing

\begin{tabbing}
\qquad\=\qquad\=\qquad\=\kill
\textsc{ComposantesFortementConnexes}\((G)\)\\
\> Construire \(G^T\)\\
\> \(\mathrm{composantes}\gets\) liste vide\\
\> \(\mathrm{non\_classes}\gets V(G)\)\\
\> Tant que \(\mathrm{non\_classes}\neq\varnothing\) :\\
\>\> Choisir \(u\in\mathrm{non\_classes}\)\\
\>\> \(X\gets\operatorname{BFS}(G,u)
           \cap\operatorname{BFS}(G^T,u)\)\\
\>\> \(\mathrm{composantes}.\mathrm{append}(X)\)\\
\>\> \(\mathrm{non\_classes}\gets\mathrm{non\_classes}\setminus X\)\\
\> Retourner \(\mathrm{composantes}\)
\end{tabbing}

Pour deux sommets \(v,w\in X\), il existe un chemin de \(v\) à \(u\) et un chemin de \(u\) à \(w\) dans \(G\).
Tous les sommets de ces chemins sont eux aussi mutuellement accessibles avec \(u\), donc appartiennent à \(X\). On peut ainsi aller de \(v\) à \(w\) entièrement dans \(G[X]\). Par conséquent, \(G[X]\) est fortement connexe.\\

De plus, si un ensemble non vide \(Z\subseteq V\setminus X\) pouvait être ajouté à \(X\) en conservant la forte connexité, chaque sommet de \(Z\) serait accessible depuis \(u\) et pourrait atteindre \(u\). Ces sommets appartiendraient donc déjà à \(X\), ce qui est une contradiction. Ainsi, \(X\) est maximal.\\

Chaque itération trouve donc une composante fortement connexe complète. Comme on choisit ensuite un sommet non classé, aucune composante n'est retournée deux fois. À la fin, tous les sommets sont classés, donc toutes les composantes
ont été trouvées.

\medskip
\textbf{Complexité.}
Posons \(n=|V|\) et \(m=|E|\). Comme discuter en (c), avec des listes d'adjacence, la construction de \(G^T\) coûte \(O(n+m)\). Chaque itération effectue un BFS dans \(G\) et un BFS dans \(G^T\), chacun de coût \(O(n+m)\). L'intersection des deux ensembles de \(X\) et \(X^t\) des sommets visités coûte \(O(n)\) : on parcourt les \(n\) sommets et on conserve ceux qui sont marqués comme visités dans les deux tableaux.

Puisque \(u\in X\), chaque itération classe au moins un nouveau sommet. Il y a donc au plus \(n\) itérations, pour une complexité totale
\[
\boxed{O\bigl(n(n+m)\bigr)}.
\]
L'algorithme est donc polynomial!
\end{enumerate}
    
\end{question}


%======================================================
%question 2
%======================================================
\begin{question}{Le problème de la sous-somme}

\begin{enumerate}[label=(\alph*)]

\item 
\end{enumerate}
\end{question}
%======================================================
%question 3
%======================================================
\begin{question}{Puissances en temps logarithmique}
    \begin{enumerate}[label=(\alph*)]
\item
On donne un contre-exemple sur un graphe complet à quatre sommets.
Les coûts des arêtes sont représentés ci-dessous.

\begin{center}
\begin{tikzpicture}[
    sommet/.style={circle, draw, fill=white, minimum size=7mm},
    cout/.style={fill=white, inner sep=2pt}
]
\node[sommet] (a) at (0,3) {$a$};
\node[sommet] (b) at (3,3) {$b$};
\node[sommet] (c) at (3,0) {$c$};
\node[sommet] (d) at (0,0) {$d$};

\draw (a) -- node[cout, above] {$1$} (b);
\draw (b) -- node[cout, right] {$1$} (c);
\draw (c) -- node[cout, below] {$10$} (d);
\draw (d) -- node[cout, left] {$3$} (a);
\draw (a) -- node[cout, pos=0.3] {$2$} (c);
\draw (b) -- node[cout, pos=0.3] {$2$} (d);
\end{tikzpicture}
\end{center}

Supposons que l'algorithme commence au sommet \(a\).
Il choisit d'abord \(b\), car \(f(ab)=1\) est le coût minimal
parmi les arêtes reliant \(a\) aux autres sommets.
Depuis \(b\), il choisit ensuite \(c\), puisque
\(f(bc)=1<f(bd)=2\).
Il ne reste alors que \(d\), puis l'algorithme revient à \(a\).
Le cycle retourné est donc
\[
C_{\mathrm{glouton}}=(a,b,c,d,a),
\qquad
f(C_{\mathrm{glouton}})=1+1+10+3=15.
\]

Pourtant, le cycle \(C'=(a,c,b,d,a)\) a un coût inférieur :
\[
f(C')=2+1+2+3=8<15.
\]

\begin{center}
\begin{tikzpicture}[
    sommet/.style={circle, draw, fill=white, minimum size=7mm},
    cout/.style={fill=white, inner sep=2pt}
]
% Cycle glouton
\begin{scope}
\draw[gray!40] (0,3) -- (3,0);
\draw[gray!40] (3,3) -- (0,0);

\draw[red, very thick] (0,3) -- node[cout, above] {$1$} (3,3);
\draw[red, very thick] (3,3) -- node[cout, right] {$1$} (3,0);
\draw[red, very thick] (3,0) -- node[cout, below] {$10$} (0,0);
\draw[red, very thick] (0,0) -- node[cout, left] {$3$} (0,3);

\node[sommet] at (0,3) {$a$};
\node[sommet] at (3,3) {$b$};
\node[sommet] at (3,0) {$c$};
\node[sommet] at (0,0) {$d$};
\node at (1.5,-0.7) {Cycle glouton : coût \(15\)};
\end{scope}

% Cycle moins coûteux
\begin{scope}[xshift=6cm]
\draw[gray!40] (0,3) -- (3,3);
\draw[gray!40] (3,0) -- (0,0);

\draw[blue, very thick] (0,3)
    -- node[cout, pos=0.3] {$2$} (3,0);
\draw[blue, very thick] (3,0)
    -- node[cout, right] {$1$} (3,3);
\draw[blue, very thick] (3,3)
    -- node[cout, pos=0.3] {$2$} (0,0);
\draw[blue, very thick] (0,0)
    -- node[cout, left] {$3$} (0,3);

\node[sommet] at (0,3) {$a$};
\node[sommet] at (3,3) {$b$};
\node[sommet] at (3,0) {$c$};
\node[sommet] at (0,0) {$d$};
\node at (1.5,-0.7) {Autre cycle : coût \(8\)};
\end{scope}
\end{tikzpicture}
\end{center}

L'algorithme peut donc retourner un cycle qui n'est pas
de coût minimum. Il ne résout pas correctement le problème
du commis voyageur : ses choix localement les moins coûteux
ne garantissent pas un coût total minimal.
\end{enumerate}
\end{question}


\bibliographystyle{apsrev4-2}
\bibliography{references}
\end{document}
